\makeatletter
\@ifundefined{isChecklistMainFile}{
  \newif\ifreproStandalone
  \reproStandalonetrue
}{
  \newif\ifreproStandalone
  \reproStandalonefalse
}
\makeatother

\ifreproStandalone
\documentclass[letterpaper]{article}
\usepackage[submission]{AuthorKit27/aaai2027}
\setlength{\pdfpagewidth}{8.5in}
\setlength{\pdfpageheight}{11in}
\frenchspacing
\begin{document}
\fi
\setlength{\leftmargini}{20pt}
\makeatletter\def\@listi{\leftmargin\leftmargini \topsep .5em \parsep .5em \itemsep .5em}
\def\@listii{\leftmargin\leftmarginii \labelwidth\leftmarginii \advance\labelwidth-\labelsep \topsep .4em \parsep .4em \itemsep .4em}
\def\@listiii{\leftmargin\leftmarginiii \labelwidth\leftmarginiii \advance\labelwidth-\labelsep \topsep .4em \parsep .4em \itemsep .4em}\makeatother

\setcounter{secnumdepth}{0}
\renewcommand\thesubsection{\arabic{subsection}}
\renewcommand\labelenumi{\thesubsection.\arabic{enumi}}

\newcounter{checksubsection}
\newcounter{checkitem}[checksubsection]

\newcommand{\checksubsection}[1]{%
  \refstepcounter{checksubsection}%
  \paragraph{\arabic{checksubsection}. #1}%
  \setcounter{checkitem}{0}%
}

\newcommand{\checkitem}{%
  \refstepcounter{checkitem}%
  \item[\arabic{checksubsection}.\arabic{checkitem}.]%
}
\newcommand{\question}[2]{\normalcolor\checkitem #1 #2 \color{blue}}
\newcommand{\ifyespoints}[1]{\makebox[0pt][l]{\hspace{-15pt}\normalcolor #1}}

\section*{Reproducibility Checklist}

\vspace{1em}
\hrule
\vspace{1em}

This checklist is filled for the current BGR submission package. The main paper is anonymous; commands, configs, per-seed CSVs, and run logs are listed in the repository's \texttt{README.md} and \texttt{results/README.md}.

\vspace{1em}
\hrule
\vspace{1em}

\checksubsection{General Paper Structure}
\begin{itemize}

\question{Includes a conceptual outline and/or pseudocode description of AI methods introduced}{(yes/partial/no/NA)}
yes. The method section defines replayable states, recovery curves, critical radii, BGR priority, and includes a training-loop figure.

\question{Clearly delineates statements that are opinions, hypothesis, and speculation from objective facts and results}{(yes/no)}
yes. The paper marks OpenVLA recovery-curve analysis as fixed-policy evidence, BGR-Suffix fine-tuning as future work, and tabular grid-policy failures as negative diagnostics.

\question{Provides well-marked pedagogical references for less-familiar readers to gain background necessary to replicate the paper}{(yes/no)}
yes. The related-work section cites curriculum learning, PLR, PAIRED/ACCEL, domain randomization, MiniGrid, LIBERO, and OpenVLA.

\end{itemize}
\checksubsection{Theoretical Contributions}
\begin{itemize}

\question{Does this paper make theoretical contributions?}{(yes/no)}
no. The paper introduces definitions and an algorithmic replay principle, but does not claim new theorems.

	\ifyespoints{\vspace{1.2em}If yes, please address the following points:}
        \begin{itemize}

	\question{All assumptions and restrictions are stated clearly and formally}{(yes/partial/no)}
	NA.

	\question{All novel claims are stated formally (e.g., in theorem statements)}{(yes/partial/no)}
	NA.

	\question{Proofs of all novel claims are included}{(yes/partial/no)}
	NA.

	\question{Proof sketches or intuitions are given for complex and/or novel results}{(yes/partial/no)}
	NA.

	\question{Appropriate citations to theoretical tools used are given}{(yes/partial/no)}
	NA.

	\question{All theoretical claims are demonstrated empirically to hold}{(yes/partial/no/NA)}
	NA.

	\question{All experimental code used to eliminate or disprove claims is included}{(yes/no/NA)}
	NA.

	\end{itemize}
\end{itemize}

\checksubsection{Dataset Usage}
\begin{itemize}

\question{Does this paper rely on one or more datasets?}{(yes/no)}
partial. Main BGR training results use procedural synthetic/grid/suffix generators. The LIBERO/OpenVLA recovery audit uses existing benchmark task definitions, init states, and fixed-policy rollout artifacts to validate learned-policy recovery-curve measurement.

\ifyespoints{If yes, please address the following points:}
\begin{itemize}

	\question{A motivation is given for why the experiments are conducted on the selected datasets}{(yes/partial/no/NA)}
	yes. The paper motivates synthetic recovery curves, procedural grid-margin replay states, robot suffix diagnostics, and LIBERO/OpenVLA as a realistic manipulation-policy target.

	\question{All novel datasets introduced in this paper are included in a data appendix}{(yes/partial/no/NA)}
	yes. The experiments are generated by checked-in scripts/configs; per-seed result CSVs are included under \texttt{results/}.

	\question{All novel datasets introduced in this paper will be made publicly available upon publication of the paper with a license that allows free usage for research purposes}{(yes/partial/no/NA)}
	yes. The generated procedural data and code are included in the repository under an MIT license.

	\question{All datasets drawn from the existing literature (potentially including authors' own previously published work) are accompanied by appropriate citations}{(yes/no/NA)}
	yes. LIBERO and OpenVLA are cited; MiniGrid is cited for grid-style context.

	\question{All datasets drawn from the existing literature (potentially including authors' own previously published work) are publicly available}{(yes/partial/no/NA)}
	yes for LIBERO task definitions, benchmark assets, and OpenVLA model code/weights. The checked-in repository contains compact derived recovery summaries; raw fixed-policy rollout artifacts are referenced by cluster path.

	\question{All datasets that are not publicly available are described in detail, with explanation why publicly available alternatives are not scientifically satisficing}{(yes/partial/no/NA)}
	NA. No non-public dataset is used for claimed quantitative results.

\end{itemize}
\end{itemize}

\checksubsection{Computational Experiments}
\begin{itemize}

\question{Does this paper include computational experiments?}{(yes/no)}
yes.

\ifyespoints{If yes, please address the following points:}
\begin{itemize}

	\question{This paper states the number and range of values tried per (hyper-) parameter during development of the paper, along with the criterion used for selecting the final parameter setting}{(yes/partial/no/NA)}
	partial. Final configs and ablation variants are checked in; the full search history is not exhaustively enumerated in the main paper.

	\question{Any code required for pre-processing data is included in the appendix}{(yes/partial/no)}
	yes. Procedural generation, LIBERO probing, OpenVLA recovery summarization, aggregation, and figure scripts are included.

	\question{All source code required for conducting and analyzing the experiments is included in a code appendix}{(yes/partial/no)}
	yes. Experiment scripts, configs, tests, aggregation, and result CSVs are checked in.

	\question{All source code required for conducting and analyzing the experiments will be made publicly available upon publication of the paper with a license that allows free usage for research purposes}{(yes/partial/no)}
	yes. The repository is on GitHub and includes an MIT license.

	\question{All source code implementing new methods have comments detailing the implementation, with references to the paper where each step comes from}{(yes/partial/no)}
	partial. Code is modular and tested; comments are intentionally sparse and do not yet cross-reference paper sections.

	\question{If an algorithm depends on randomness, then the method used for setting seeds is described in a way sufficient to allow replication of results}{(yes/partial/no/NA)}
	yes. Every config lists seeds; scripts pass seeds to NumPy generators and result CSVs include per-seed rows.

	\question{This paper specifies the computing infrastructure used for running experiments (hardware and software), including GPU/CPU models; amount of memory; operating system; names and versions of relevant software libraries and frameworks}{(yes/partial/no)}
	yes. The run ledger records Slurm partitions, CPU/memory/time requests, hostnames for logged runs, environment constraints, and JSON environment snapshots for compute and GPU jobs.

	\question{This paper formally describes evaluation metrics used and explains the motivation for choosing these metrics}{(yes/partial/no)}
	yes. Recovery AUC, critical radius, clean success, transfer RAUC, and AULC are defined or described.

	\question{This paper states the number of algorithm runs used to compute each reported result}{(yes/no)}
	yes. Tables and text state three- or five-seed means as applicable.

	\question{Analysis of experiments goes beyond single-dimensional summaries of performance (e.g., average; median) to include measures of variation, confidence, or other distributional information}{(yes/no)}
	yes. Aggregated tables include standard errors; experiments report median critical radius and AULC in addition to final means.

	\question{The significance of any improvement or decrease in performance is judged using appropriate statistical tests (e.g., Wilcoxon signed-rank)}{(yes/partial/no)}
	yes. The repository includes paired exact sign-flip tests for the reported seed-paired comparisons. Because the current diagnostic runs use only three or five seeds, the paper interprets these tests descriptively and does not overstate thresholded significance.

	\question{This paper lists all final (hyper-)parameters used for each model/algorithm in the paper's experiments}{(yes/partial/no/NA)}
	yes. Final hyperparameters are in the checked-in YAML configs under \texttt{configs/}.

\end{itemize}
\end{itemize}
\ifreproStandalone
\end{document}
\fi
